% Backup frames for the seminar "AI in the Flipped Classroom" (Fellowship in Education, 31 Oct 2026).
% Used after the thanks slide by Main_Seminar_AI_For_Educators_Flipped_Presentation.tex (inside \appendix,
% so they are not in the frame count) and at the end of the _CheatSheet.tex Prompt Pack.
% Moved out of the live deck on 2026-10-26 at the author's request; the frames are unchanged.
%%%%%%%%%%%%%%%%%%%%%%%%%%%%%%%%%%%%%%%%%%%%%%%%%%%
\begin{frame}[fragile]\frametitle{}
\begin{center}
{\Large Backup: Fellowship Project Ideas}

\medskip
{\small (Use if the group wants project work; also the content-routes table)}
\end{center}
\end{frame}

%%%%%%%%%%%%%%%%%%%%%%%%%%%%%%%%%%%%%%%%%%%%%%%%%%%
\begin{frame}[fragile]\frametitle{Backup: Choosing How Students Get the Content}
\begin{center}
\resizebox{\linewidth}{!}{%
\begin{tabular}{lcccc}
\toprule
 & Textbook & Own notes & Video playlist & LMS or MOOC \\
\midrule
Teacher effort & Low & Medium & Low & High \\
Matches your syllabus & Partly & Fully & Partly & Fully \\
Interactive & No & No & Little & Yes \\
Tracks students & No & No & No & Yes \\
\bottomrule
\end{tabular}}
\end{center}
\begin{itemize}
\item Routes combine well: one LMS unit can hold your notes, a video and a quiz.
\item Moodle and Open edX are two free LMS platforms. Hosting takes effort, so today uses plain notes.
\end{itemize}
\end{frame}

%%%%%%%%%%%%%%%%%%%%%%%%%%%%%%%%%%%%%%%%%%%%%%%%%%%
\begin{frame}[fragile]\frametitle{Exercise: Fellowship Project Ideas}
\begin{itemize}
\item \textbf{Goal:} three topic options and one draft problem statement for your fellowship project.
\item \textbf{Step 1:} describe your real setting to the AI and ask for topic options.
\item \textbf{Step 2:} answer its questions, and choose one topic.
\item \textbf{Step 3:} write the problem statement yourself, then ask the AI to critique it.
\item \textbf{Result:} a topic you can take into the proposal-writing session.
\end{itemize}
\end{frame}

%%%%%%%%%%%%%%%%%%%%%%%%%%%%%%%%%%%%%%%%%%%%%%%%%%%
\begin{frame}[fragile]\frametitle{Exercise: Fellowship Project Ideas, Topic Prompt}
\begin{itemize}
\item Be honest about your constraints: devices, internet, class size, time.
\end{itemize}
\begin{lstlisting}[language={}]
You are an education researcher mentoring a school
teacher (R). I teach [SUBJECT] to [GRADE] in a
[TYPE OF SCHOOL]. My constraints: [DEVICES, INTERNET,
CLASS SIZE, TIME] (C).
Suggest 5 project topics on using a flipped
classroom with AI in my setting. For each: the
problem it solves and one measurable outcome (A).
Then ask me 5 questions to choose the best one (T).
\end{lstlisting}
\end{frame}

%%%%%%%%%%%%%%%%%%%%%%%%%%%%%%%%%%%%%%%%%%%%%%%%%%%
\begin{frame}[fragile]\frametitle{Exercise: Fellowship Project Ideas, Critique Prompt}
\begin{itemize}
\item Write 3 to 4 sentences yourself first. The AI critiques; it does not write your proposal.
\end{itemize}
\begin{lstlisting}[language={}]
Here is my draft problem statement:
[PASTE YOUR 3 TO 4 SENTENCES]
List its weaknesses: vague words, no measure,
too big for one year, missing baseline.
Then list 5 questions a reviewer would ask me.
Do not rewrite it.
\end{lstlisting}
\end{frame}

%%%%%%%%%%%%%%%%%%%%%%%%%%%%%%%%%%%%%%%%%%%%%%%%%%%
\begin{frame}[fragile]\frametitle{Exercise: Fellowship Project Ideas, Check}
\begin{itemize}
\item \textbf{Your class:} does the topic fit your real students, school and year?
\item \textbf{Measurable:} could you collect the outcome with tools you have?
\item \textbf{References:} the AI may invent papers and links. Search for each one yourself.
\item \textbf{Ownership:} the proposal text is yours. Use the AI as a mentor.
\end{itemize}
\end{frame}
